\section*{第 \ref{ch:g8:ohms-law} 章——电阻与欧姆定律}

\begin{solution}{exo:g8:ohms-law:1}
\omterm{def:g8:ohms-law:resistance}{电阻}是元件对电流阻碍作用的强弱；\omterm{def:g6:measurement-in-science:measuring}{单位}为\omterm{def:g8:ohms-law:resistance}{欧姆}，
符号为 $\Omega$。$U = R I$；\ $I = U/R$；\ $R = U/I$。
\end{solution}

\begin{solution}{exo:g8:ohms-law:2}
$I = 6 \div 30 = \qty{0.2}{A}$; \ $U = 30 \times 0.5 =
\qty{15}{V}$.
\end{solution}

\begin{solution}{exo:g8:ohms-law:3}
$R = 4.5 \div 0.09 = 50\,\Omega$.
\end{solution}

\begin{solution}{exo:g8:ohms-law:4}
过原点的直线——正比例关系。
斜率越大的直线，表示每\omterm{def:g8:intensity-ammeter:intensity}{安培}需要的\omterm{def:g8:voltage-voltmeter:voltage}{电压}越高，
即\omterm{def:g8:ohms-law:resistance}{电阻}越大。
\end{solution}

\begin{solution}{exo:g8:ohms-law:5}
$3.0 \div 0.20 = 15$；$6.0 \div 0.40 = 15$。各行比值相同，
\emph{正是}\omterm{def:g8:ohms-law:resistance}{欧姆}定律：一个常数 $R$ 对应
每一次读数，使“\omterm{def:g8:ohms-law:resistance}{电阻}”成为
元件本身的性质。
\end{solution}

\begin{solution}{exo:g8:ohms-law:6}
\omterm{def:g8:ohms-law:resistance}{电阻}越小，取用的电流越大：$230 \div 25 =
\qty{9.2}{A}$，对比 $230 \div 35 \approx \qty{6.6}{A}$。
\end{solution}

\begin{solution}{exo:g8:ohms-law:7}
\omterm{def:g4:series-circuits:series}{串联}的\omterm{def:g3:first-electric-circuit:bulb}{灯泡}共用同一电流 $I$；每只\omterm{def:g3:first-electric-circuit:bulb}{灯泡}的\omterm{def:g8:voltage-voltmeter:voltage}{电压}为 $U = R I$，
电流 $I$ 相同——故\omterm{def:g8:ohms-law:resistance}{电阻} $R$ 越大，
对应的\omterm{def:g8:voltage-voltmeter:voltage}{电压} $U$ 也越大：阻碍作用更强的元件分得的\omterm{def:g8:voltage-voltmeter:voltage}{电压}更多。
\end{solution}

\begin{solution}{exo:g8:ohms-law:8}
\omterm{def:g7:short-circuits-safety:device}{短路}：$R$ 接近零，故 $I = U/R$ 急剧增大——电流狂涌
就是除以了几乎为零的数。断开的\omterm{ex:g3:first-electric-circuit:switch}{开关}：$R$
无限大，故 $I = U/R$ 降为零——
完全阻断。
\end{solution}

\begin{solution}{exo:g8:ohms-law:9}
\omterm{def:g8:ohms-law:resistor}{电阻器}须分得 $9 - 2 = \qty{7}{V}$ 的\omterm{def:g8:voltage-voltmeter:voltage}{电压}，电流为 \qty{0.02}{A}：
$R = 7 \div 0.02 = 350\,\Omega$.
\end{solution}

\begin{solution}{exo:g8:ohms-law:10}
第一个点：$R = 1 \div 0.25 = 4\,\Omega$；第二个点：$4 \div
0.5 = 8\,\Omega$——没有单一的值与之相符。其原因在于：
电流增大使灯丝发热，而热的金属
\omterm{def:g8:ohms-law:resistance}{电阻}更大；图线之所以弯曲，
是因为元件在工作过程中自身发生了变化。
\end{solution}

\begin{solution}{exo:g8:ohms-law:11}
同一回路中只有一个电流 $I$：两段\omterm{def:g8:voltage-voltmeter:voltage}{电压}降之和等于\omterm{def:g3:first-electric-circuit:battery}{电池}\omterm{def:g8:voltage-voltmeter:voltage}{电压}，
$6 = 10 I + 20 I = 30 I$，故 $I = \qty{0.2}{A}$。随后 $U_{10}
= 10 \times 0.2 = \qty{2}{V}$ 和 $U_{20} = 20 \times 0.2 =
\qty{4}{V}$——二者之和为 \qty{6}{V}，符合
回路\omterm{def:g8:voltage-voltmeter:voltage}{电压}定律。
\end{solution}

\begin{solution}{exo:g8:ohms-law:12}
每条\omterm{def:g5:series-parallel-circuits:parallel}{支路}两端\omterm{def:g8:voltage-voltmeter:voltage}{电压}均为 \qty{6}{V}：$I_{10} = 6 \div 10 =
\qty{0.6}{A}$，$I_{20} = 6 \div 20 = \qty{0.3}{A}$；总电流
$\qty{0.9}{A}$。等效\omterm{def:g8:ohms-law:resistance}{电阻}：$R = 6 \div 0.9 \approx
6.7\,\Omega$——\emph{小于}其中任一\omterm{def:g8:ohms-law:resistance}{电阻}。
用道路作比喻：连接同样两地的两条道路
承载的车流多于任何一条单独的道路；每增加一条\omterm{def:g5:series-parallel-circuits:parallel}{并联}道路，
都会使通行更顺畅，因此整体的阻碍
低于其中最弱一员的阻碍。
\end{solution}

\begin{solution}{pb:g8:ohms-law:1}
\textbf{1.} 电源、\omterm{def:g8:intensity-ammeter:ammeter}{电流表}与元件\omterm{def:g4:series-circuits:series}{串联}成一个回路，
好比道路上的收费站；\omterm{def:g8:voltage-voltmeter:voltmeter}{电压表}跨接在
元件两端——好比路边的\omterm{def:g6:measurement-in-science:measuring}{测量}员。
\textbf{2.} $R = 4.5 \div 0.15 = 30\,\Omega$.
\textbf{3.} $33$ 的百分之五约为 $1.65$，可接受的范围是
$31.35$ 至 $34.65\,\Omega$ 之间。样品一的 $30\,\Omega$
不在此范围内：应予剔除。
\textbf{4.} $I = 4.5 \div 150 = \qty{0.03}{A} =
\qty{30.0}{mA}$.
\textbf{5.} 每行都得出 $R = U/I = 30\,\Omega$（$1.5
\div 0.05$，$3.0 \div 0.10$，\dots）：比值恒定，
图线为直线——一个数值适用于全部读数：$30\,\Omega$。
\textbf{6.} 各比值依次为 $5$、$6.7$、$8.3$、$10\,\Omega$：
逐步攀升——\omterm{def:g8:ohms-law:resistance}{电阻}随电流增大而增大；
不存在单一的 $R$。
\textbf{7.} 电流增大时，灯丝\omterm{def:g3:temperature-thermometers:temperature}{温度}升高，
直至发光，热金属的\omterm{def:g8:ohms-law:resistance}{电阻}更大——
图线向上弯曲。因此，这一直线测试恰好只接受
$R$ 保持不变的元件：按其构造，它
就是一台\omterm{def:g8:ohms-law:resistor}{电阻器}检测器。
\textbf{8.} 稳定\omterm{def:g3:temperature-thermometers:temperature}{温度}条款：\omterm{def:g8:ohms-law:resistance}{欧姆}定律所保证的正比关系
\emph{仅在\omterm{def:g3:temperature-thermometers:temperature}{温度}稳定时成立}——若在\omterm{def:g6:measurement-in-science:measuring}{测量}途中加热元件，
即使是合格的\omterm{def:g8:ohms-law:resistor}{电阻器}，其读数也会漂移。
\textbf{9.} $R = 12 \div 0.03 = 400\,\Omega$.
\textbf{10.} $U = R I = 25 \times 9.2 = \qty{230}{V}$ ---
恰好是市电：线圈取用的，正是\omterm{def:g8:ohms-law:resistance}{欧姆}定律
分配给它的数值，\omterm{def:g7:short-circuits-safety:fuse}{保险丝}的余量也按设计运作。
没有什么“只是”的东西，一切都是算术。
\textbf{11.} 跨接线生来就是一根连接导线——
平坦的道路，不消耗任何\omterm{def:g8:voltage-voltmeter:voltage}{电压}。误接在电源两端时，
其近乎为零的\omterm{def:g8:ohms-law:resistance}{电阻}使它成为完美的\omterm{def:g7:short-circuits-safety:device}{短路}：这个“反派”
其实始终只是放错位置的一根好导线。
\textbf{12.} 例如：“定律：$U = R I$，\omterm{def:g8:voltage-voltmeter:voltage}{电压}等于
阻碍乘以电流。其图线为过原点的直线，
阻碍越强，斜率越大。唯有
\omterm{def:g3:temperature-thermometers:temperature}{温度}稳定的元件才有资格为其背书——\omterm{def:g5:heat-and-insulation:heat}{热量}会改写
\omterm{def:g8:ohms-law:resistance}{电阻}，发热时绘制的图线毫无说服力。”
\end{solution}
